% =============================================================================
%  lec08.tex —— 第 8 讲 极限与余极限：模型范畴（底本 PDF 第 66–72 页）
% =============================================================================
\lecture{8}{极限与余极限：模型范畴}
\begin{center}\small\color{black!55}（底本 PDF 第 66–72 页）\end{center}

本讲用模型范畴来计算极限与余极限。

\section{$\infty$-极限作为同伦极限}

\nitem\label{L8-1.1} 回忆我们有 Quillen 等价
\[
\fC\colon \Set_\Delta^{\mathrm{Joyal}}\rightleftarrows \Cat_\Delta : N_\bullet .
\]
本节将解释下述结果：

\begin{theorem}[HTT.4.2.4.1]\label{L8-1.2}
设 $\cC$ 是组合单纯模型范畴。则对小的纤维单纯范畴 $\cJ$，图
\[
\cJ^{\triangleleft}\to \cC^{\circ}
\]
是 $\cC$ 中的同伦极限图，当且仅当相应的图
\[
N_\bullet(\cJ)^{\triangleleft}\to N_\bullet(\cC^{\circ})
\]
是拟范畴 $N_\bullet(\cC^{\circ})$ 中的极限图。
\end{theorem}

\begin{variant}\label{L8-1.3}
对偶地，图 $\cJ^{\triangleright}\to \cC^{\circ}$ 是同伦余极限图，当且仅当 $N_\bullet(\cJ)^{\triangleright}\to N_\bullet(\cC^{\circ})$ 是极限图。
\end{variant}

\yizhu{原文此句末作“is a limit diagram”。按对偶，同伦余极限图所对应的应是 $N_\bullet(\cC^{\circ})$ 中的余极限图，此处疑为笔误；译文照录原文。\textbf{[来源：A]}（HTT.4.2.4.1 的对偶陈述）}

\nitem\label{L8-1.4} 我们很快就会给出定理~\ref{L8-1.2} 中未定义概念的准确定义。眼下，先满足于下面这些非正式的说法。

\textbf{单纯模型范畴}是配有相容单纯充实的模型范畴 $\cC$。这里的相容性条件保证：
\begin{enumerate}[label=(\arabic*),leftmargin=2.4em,itemsep=0.2em]
\item 有典范等价
\[
h\cC\simeq \pi_0\cC,
\]
其中 $h\cC$ 是模型范畴 $\cC$ 的同伦范畴（见第 2 讲定义 2.20），而 $\pi_0\cC$ 是单纯范畴 $\cC$ 的同伦范畴（见第 5 讲定义 5.3）。
\item 对双纤维性的对象 $x,y\in\cC$，单纯集 $\Hom_{\cC}(x,y)$ 是 Kan 复形。换言之，$\cC^{\circ}$ 是 $\Cat_\Delta$ 中的纤维对象。
\end{enumerate}

对模型范畴而言，组合性是一个技术性的集合论大小条件，眼下可以忽略。

对单纯模型范畴 $\cC$，图 $u\colon \cJ\to\cC$ 的\textbf{同伦极限}是朴素极限函子的右导出函子的取值。换言之，
\[
\holim u\coloneqq \R\lim(u),
\]
其中函子 $\lim$ 是 Quillen 伴随
\[
\mathrm{const}\colon \cC\rightleftarrows \Fun(\cJ,\cC) : \lim
\]
中的右伴随。这里 $\Fun(\cJ,\cC)$ 是从 $\cJ$ 到 $\cC$ 的单纯充实函子范畴，配有由 $\cC$ 上模型结构诱导的适当模型结构。

\nitem\label{L8-1.5} 回忆 $\Set_\Delta$ 是笛卡尔闭的，因此带有由 $\Fun(-,-)$ 给出的自然单纯充实。这个充实与 Kan--Quillen 模型结构相容。我们把所得的单纯模型范畴记为 $\Set^{\mathrm{KQ}}_\Delta$。在定理~\ref{L8-1.2} 中可取 $\cC\coloneqq \Set^{\mathrm{KQ}}_\Delta$。注意 $\cC^{\circ}=\Kan$ 就是第 5 讲 §7 中定义的单纯范畴。于是得到：

\begin{corollary}\label{L8-1.6}
设 $\cJ$ 是小的纤维单纯范畴。则 $\Kan\subset \Set^{\mathrm{KQ}}_\Delta$ 中 $\cJ$-指标图的同伦（余）极限，计算拟范畴
\[
\Kan \coloneqq N_\bullet(\Kan)
\]
中相应的 $N_\bullet(\cJ)$-指标图的（余）极限。
\end{corollary}

\begin{remark}\label{L8-1.7}
注意后者建模了 $\infty$-范畴 $\Grpd_\infty$ 中的（余）极限，而它们本质上控制任意 $\infty$-范畴中的（余）极限（第 7 讲定理 2.11）。
\end{remark}

\begin{remark}\label{L8-1.8}
上述推论给出了计算 $\infty$-群胚的小（余）极限的模型范畴算法。我们以第二种情形为例。设 $u\colon K\to \Grpd_\infty$ 是 $K\in\Set_\Delta$ 上的小图。
\begin{enumerate}[label=(\roman*),leftmargin=2.8em,itemsep=0.2em]
\item 选取 $\Cat_\Delta$ 中的弱等价 $\fC(K)\to \cJ$。
\item 找函子 $w\colon \cJ\to \Kan$，使得复合
\[
K\to N_\bullet(\cJ)\to N_\bullet(\Kan)\eqqcolon \Kan
\]
表示 $u$。由 HTT.4.2.4.4，这样的 $w$ 总存在。
\item 把 $w$ 视为模型范畴 $\Fun(\cJ,\Set^{\mathrm{KQ}}_\Delta)$ 中的对象。通过求 $w$ 的（上）纤维替换，计算导出（余）极限
\[
\R\lim(w),\ \mathbb{L}\colim(w)\in \Set^{\mathrm{KQ}}_\Delta[W^{-1}].
\]
\end{enumerate}
由推论~\ref{L8-1.6}，所得的 $\Set^{\mathrm{KQ}}_\Delta[W^{-1}]\simeq \hGrpd_\infty$ 中的对象典范同构于 $\lim u$ 与 $\colim u$。
\end{remark}

\begin{exercise}\label{L8-1.9}
证明：当 $\fC(K)\to \cJ$ 是上纤维化时，步骤 (ii) 中的函子 $w$ 存在。
\end{exercise}

\begin{remark}\label{L8-1.10}
在注~\ref{L8-1.8} 中，当 $K=N_\bullet(J)$ 是普通范畴的脉时，可以取 $\cJ\coloneqq J$，即把 $J$ 视为带离散充实的单纯范畴。此时 $\Fun(J,\Set^{\mathrm{KQ}}_\Delta)$ 就是普通范畴之间的函子范畴 $J\to\Set_\Delta$，而未导出的函子 $\lim$ 就是普通范畴的极限函子。

事实上，上述讨论本质上涵盖了所有情形：因为对任意单纯集 $K$，都存在初始态射 $N_\bullet(J)\to K$，其中 $J$ 是普通范畴，甚至可以是偏序集。见 HTT.4.2.3.14。
\end{remark}

\begin{remark}\label{L8-1.11}
在注~\ref{L8-1.8} 中，所得对象落在同伦范畴 $\hGrpd_\infty$ 中，而不是 $\Grpd_\infty$ 中。换言之，上述算法只把 $\lim u$ 计算到同伦为止。特别地，它既不能给出整个极限图，也不能给出 $\lim u$ 在 $(\Grpd_\infty)_{/u}$ 中的典范提升。

为补救这一点，可以尝试找扩张图 $\overline{w}\colon \cJ^{\triangleleft}\to \Kan$，它把 $\overline{w}(\ast)$ 展示为 $w$ 的同伦极限（见下文定义~\ref{L8-2.18}）。于是定理~\ref{L8-1.2} 说，相应的图
\[
N_\bullet(\cJ)^{\triangleleft}\to N(\Kan)\eqqcolon \Kan
\]
也是极限图。现在，沿范畴等价 $K^{\triangleleft}\to N_\bullet(\cJ)^{\triangleleft}$ 作限制，就得到极限图，它扩张 $v\colon K\to\Kan$。

对一般的 $w$，上述扩张 $\overline{w}$ 可能不存在\footnote{例如，考虑 $\Sing_\bullet(\{0\}\to[0,1]\leftarrow\{1\})$。}。

然而，我们可以把 $w$ 换成它的任意纤维替换，因为它们在 $\Grpd_\infty$ 中表示同一个图。在这个附加假设下，这样的扩张 $\overline{w}$ 存在，并且在唯一等价的意义下唯一，因为它必须是扩张 $w$ 的极限图。
\end{remark}

\begin{remark}\label{L8-1.12}
在以后的讲次中我们会看到，任意可表现 $\infty$-范畴都可实现为某个组合单纯模型范畴 $\cC$ 的 $N_\bullet(\cC^{\circ})$。这给出了计算任意可表现 $\infty$-范畴中小（余）极限的模型范畴算法。
\end{remark}

\section{同伦极限的定义}

\nitem\label{L8-2.1} 本节给出定理~\ref{L8-1.2} 中所用概念的准确定义。

\begin{definition}\label{L8-2.2}
设 $\cC$ 是单纯范畴。
\begin{enumerate}[label=(\arabic*),leftmargin=2.4em,itemsep=0.2em]
\item 若对任意 $S\in\Set_\Delta$ 与 $X\in\cC$，函子
\[
\Fun(S,\Hom_{\cC}(X,-))\colon \cC\to \Set_\Delta
\]
可由 $\cC$ 中的某个对象表示，则称 $\cC$ 在 $\Set_\Delta$ 上\textbf{张量}（tensored），并把该表示对象记为 $S\otimes X$。
\item 若对任意 $S\in\Set_\Delta$ 与 $Y\in\cC$，函子
\[
\Fun(S,\Hom_{\cC}(-,Y))\colon \cC^{\op}\to \Set_\Delta
\]
可由 $\cC$ 中的某个对象表示，则称 $\cC$ 在 $\Set_\Delta$ 上\textbf{余张量}（cotensored），并把该表示对象记为 $\Fun(S,Y)$。
\end{enumerate}
\end{definition}

\begin{theoremdef}[HTT.A.3.1.5, A.3.1.6]\label{L8-2.3}
设 $\cC$ 是配有单纯充实、且在 $\Set_\Delta$ 上既张量又余张量的模型范畴。则下列条件等价：
\begin{enumerate}[label=(\roman*),leftmargin=2.8em,itemsep=0.2em]
\item 给定 $\cC$ 中任意上纤维化 $j\colon X\to X'$ 与任意纤维化 $k\colon Y\to Y'$，态射
\[
\Hom_{\cC}(X',Y)\to \Hom_{\cC}(X',Y')\times_{\Hom_{\cC}(X,Y')}\Hom_{\cC}(X,Y)
\]
是 $\Set^{\mathrm{KQ}}_\Delta$ 中的纤维化，且当 $j$ 或 $k$ 是弱等价时它也是弱等价。
\item 给定 $\Set^{\mathrm{KQ}}_\Delta$ 与 $\cC$ 中分别的上纤维化 $i\colon S\to S'$ 与 $j\colon X\to X'$，态射
\[
(S'\otimes X)\sqcup_{S\otimes X}(S\otimes X')\to S'\otimes X'
\]
是 $\cC$ 中的上纤维化，且当 $i$ 或 $j$ 是弱等价时它也是弱等价。
\item 给定 $\Set^{\mathrm{KQ}}_\Delta$ 中的上纤维化 $i\colon S\to S'$ 与 $\cC$ 中的纤维化 $k\colon Y\to Y'$，态射
\[
\Fun(S',Y)\to \Fun(S',Y')\times_{\Fun(S,Y')}\Fun(S,Y)
\]
是 $\cC$ 中的纤维化，且当 $i$ 或 $k$ 是弱等价时它也是弱等价。
\end{enumerate}
若上述条件满足，则称 $\cC$ 是\textbf{单纯模型范畴}。
\end{theoremdef}

\begin{exercise}\label{L8-2.4}
设 $\cC$ 是单纯模型范畴。证明：若 $X$ 上纤维且 $Y$ 纤维，则 $\Hom_{\cC}(X,Y)$ 是 Kan 复形。
\end{exercise}

\begin{exercise}\label{L8-2.5}
你能把 $\Set^{\mathrm{Joyal}}_\Delta$ 变成单纯模型范畴吗？
\end{exercise}

\begin{exercise}\label{L8-2.6}
设 $\cC$ 是单纯模型范畴。构造典范等价 $h\cC\simeq \pi_0\cC$。
\end{exercise}

\nitem\label{L8-2.7} 我们还需要下述技术性大小条件。

\begin{definition}[HTT.A.2.6.1]\label{L8-2.8}
设 $\cC$ 是模型范畴。若下列条件满足，则称 $\cC$ 是\textbf{组合的}：
\begin{enumerate}[label=(\alph*),leftmargin=2.4em,itemsep=0.1em]
\item 范畴 $\cC$ 可表现。
\item 作为弱饱和的态射类，$(C)$ 由一个集合生成。
\item 作为弱饱和的态射类，$(C\cap W)$ 由一个集合生成。
\end{enumerate}
\end{definition}

\begin{example}\label{L8-2.9}
配有单纯充实 $\Fun(-,-)$ 的模型范畴 $\Set^{\mathrm{KQ}}_\Delta$ 是组合单纯模型范畴。
\end{example}

\begin{theoremdef}[HTT.A.3.3.2]\label{L8-2.10}
设 $\cC$ 是组合单纯模型范畴，$\cJ$ 是小单纯范畴。则 $\Fun(\cJ,\cC)$ 上有两个组合模型结构：
\begin{enumerate}[label=(\arabic*),leftmargin=2.4em,itemsep=0.3em]
\item \textbf{投射模型结构}，记为 $\Fun(\cJ,\cC)_{\mathrm{proj}}$，其中
\begin{itemize}[leftmargin=2.8em,itemsep=0.1em]
\item[(W)] 弱等价是逐点弱等价的自然变换。
\item[(F)] 投射纤维化是逐点纤维化的自然变换。
\item[(C)] 投射上纤维化的类由 $(C\cap W)$ 决定。
\end{itemize}
\item \textbf{内射模型结构}，记为 $\Fun(\cJ,\cC)_{\mathrm{inj}}$，其中
\begin{itemize}[leftmargin=2.8em,itemsep=0.1em]
\item[(W)] 弱等价是逐点弱等价的自然变换。
\item[(C)] 内射上纤维化是逐点上纤维化的自然变换。
\item[(F)] 内射纤维化的类由 $(F\cap W)$ 决定。
\end{itemize}
\end{enumerate}
\end{theoremdef}

\begin{example}\label{L8-2.11}
当 $\cJ=[0]$ 是单点时，两个模型结构都与 $\cC$ 上给定的模型结构重合。
\end{example}

\begin{proposition}[HTT.A.3.3.6]\label{L8-2.12}
设
\[
F\colon \cA\rightleftarrows \cB : G
\]
是组合单纯模型范畴之间的 Quillen 伴随。则对任意小单纯范畴 $\cJ$，它诱导 Quillen 伴随
\[
F\circ -\colon \Fun(\cJ,\cA)_{?}\rightleftarrows \Fun(\cJ,\cB)_{?} : G\circ -,
\]
其中 $?$ 可以是 proj 或 inj。若 $F\colon \cA\rightleftarrows \cB : G$ 是 Quillen 等价，则它们也是 Quillen 等价。
\end{proposition}

\begin{proposition}[HTT.A.3.3.7, A.3.3.8]\label{L8-2.13}
设 $\iota\colon \cJ\to\cJ'$ 是小单纯充实范畴之间的函子。对组合单纯模型范畴 $\cC$，我们有 Quillen 伴随
\[
\mathrm{LKE}_{\iota}\colon \Fun(\cJ,\cC)_{\mathrm{proj}}\rightleftarrows \Fun(\cJ',\cC)_{\mathrm{proj}} : \iota\circ -;
\]
\[
\iota\circ -\colon \Fun(\cJ',\cC)_{\mathrm{inj}}\rightleftarrows \Fun(\cJ,\cC)_{\mathrm{inj}} : \mathrm{RKE}_{\iota};
\]
若 $\iota\colon \cJ\to\cJ'$ 是模型范畴 $\Cat_\Delta$ 中的弱等价（见第 5 讲定义 5.4），则它们是 Quillen 等价。
\end{proposition}

\begin{definition}\label{L8-2.14}
我们把
\[
\mathrm{RKE}_{\iota}\colon \Fun(\cJ,\cC)_{\mathrm{inj}}\to \Fun(\cJ',\cC)_{\mathrm{inj}}
\]
的右导出函子称为\textbf{同伦右 Kan 扩张函子}，并记作
\[
\ho\mathrm{RKE}_{\iota}\colon \Fun(\cJ,\cC)[W^{-1}]\to \Fun(\cJ',\cC)[W^{-1}].
\]
当 $\cJ'=[0]$ 是单点时，得到\textbf{同伦极限函子}
\[
\holim\colon \Fun(\cJ,\cC)[W^{-1}]\to \cC[W^{-1}],
\]
它是
\[
\lim\colon \Fun(\cJ,\cC)_{\mathrm{inj}}\to \cC
\]
的右导出函子。

对偶地，我们定义\textbf{同伦左 Kan 扩张函子}与\textbf{同伦余极限函子}。
\end{definition}

\begin{warning}\label{L8-2.15}
注意
\[
\Fun(\cJ,\cC)[W^{-1}]\neq \Fun(\cJ,\cC[W^{-1}]).
\]
因此同伦极限并不是关于 $\cC[W^{-1}]$ 中图的函子性构造。
\end{warning}

\begin{remark}\label{L8-2.16}
事实上，定理~\ref{L8-1.2} 表明 $\infty$-极限的理论可作为经典同伦极限理论非函子性的补救。也就是说，人们不应考虑映入普通同伦范畴 $\cC[W^{-1}]$ 的图，而应考虑映入由拟范畴 $N_\bullet(\cC^{\circ})$ 建模的 $\infty$-范畴的图。事实上，后者可典范地等同于拟范畴局部化 $N_\bullet(\cC)[W^{-1}]$。更多信息见第 5 讲 A.5。
\end{remark}

\begin{construction}\label{L8-2.17}
设 $\overline{w}\colon \cJ^{\triangleleft}\to\cC$ 是函子，$w\colon \cJ\to\cC$ 是它的限制。$\Fun(\cJ,\cC)$ 中有从常值函子 $\overline{w}(\ast)$ 到 $w$ 的显然态射。由伴随性，我们得到典范态射
\[
\overline{w}(\ast)\to \lim w.
\]
\end{construction}

\begin{definition}\label{L8-2.18}
设 $\overline{w}\colon \cJ^{\triangleleft}\to\cC$ 是函子，$w\colon \cJ\to\cC$ 是它的限制。若对任意（等价地，对每个）纤维替换 $w\to w'$，复合
\[
\overline{w}(\ast)\to \lim w\to \lim w'
\]
是同构，则称 $\overline{w}$ 把 $\overline{w}(\ast)$ \textbf{展示}为 $w$ 的同伦极限。
\end{definition}

\yizhu{原文此处作“for any/all fibrant replacement”，意为：只要该条件对某个纤维替换成立，就对一切纤维替换成立，故与替换的选取无关。\textbf{[来源：I]}}

\section{同伦（余）极限的例子}

\nitem\label{L8-3.1} 本节的 $\cC$ 始终是组合单纯模型范畴。

\begin{exercise}\label{L8-3.2}
设 $\cJ$ 是集合。证明：
\begin{enumerate}[label=(\arabic*),leftmargin=2.4em,itemsep=0.1em]
\item 函子 $w\colon \cJ\to\cC$ 在 $\Fun(\cJ,\cC)_{\mathrm{inj}}$ 中纤维，当且仅当对任意 $j\in\cJ$，$w(j)\in\cC$ 纤维。
\item 函子 $w\colon \cJ\to\cC$ 在 $\Fun(\cJ,\cC)_{\mathrm{proj}}$ 中上纤维，当且仅当对任意 $j\in\cJ$，$w(j)\in\cC$ 上纤维。
\end{enumerate}
由此推出：（上）纤维对象的同伦（余）积可由朴素（余）积计算。
\end{exercise}

\begin{exercise}\label{L8-3.3}
设 $\cJ\coloneqq \{a_0\xrightarrow{f_0} b\xleftarrow{f_1} a_1\}$ 是拉回的指标范畴。证明：函子 $w\colon \cJ\to\cC$ 在 $\Fun(\cJ,\cC)_{\mathrm{inj}}$ 中纤维，当且仅当 $w(b)\in\cC$ 纤维且 $w(f_i)$ 都是纤维化。由此推出 $\cC$ 中同伦拉回图的一个充分条件。
\end{exercise}

\begin{remark}\label{L8-3.4}
注意上述条件比第 1 讲练习 A.1 中的那些条件更强。原因在于：后者是借助 $\Fun(\cJ,\cC)$ 上另一个模型结构得到的，即所谓的 Reedy 模型结构。更多信息见 HTT.A.2.9。
\end{remark}

\begin{exercise}\label{L8-3.5}
设 $\cJ\coloneqq \{a\rightrightarrows b\}$ 是等化子的指标范畴。证明：函子 $w\colon \cJ\to\cC$ 在 $\Fun(\cJ,\cC)_{\mathrm{inj}}$ 中纤维，当且仅当 $w(b)\in\cC$ 纤维且 $w(a)\to w(b\times b)$ 是纤维化。由此推出 $\cC$ 中同伦等化子图的一个充分条件。
\end{exercise}

\begin{exercise}\label{L8-3.6}
设 $\cJ\coloneqq \{\cdots<-2<-1<0\}$ 是序列极限的指标范畴。证明：函子 $w\colon \cJ\to\cC$ 在 $\Fun(\cJ,\cC)_{\mathrm{inj}}$ 中纤维，当且仅当 $w(0)\in\cC$ 纤维且 $w(-n)\to w(-n+1)$ 是纤维化。由此推出 $\cC$ 中同伦序列极限图的一个充分条件。
\end{exercise}

\begin{exercise}\label{L8-3.7}
你能为 $\cC$ 中的同伦序列极限图找到更弱的必要条件吗？
\end{exercise}

\begin{exercise}\label{L8-3.8}
设 $\cJ$ 是 $\Cat_\Delta$ 中的纤维对象，$\underline{X}\colon \cJ\to\cC$ 是取值为 $X$ 的常值函子。利用第 7 讲附录 B 的结果证明：
\[
\holim \underline{X}\simeq \Fun(N_\bullet(\cJ),X),\qquad
\hocolim \underline{X}\simeq N_\bullet(\cJ)\otimes X.
\]
\end{exercise}

\appendixsection{标记单纯集}

\begin{definition}\label{L8-A.1}
设 $\Set^+_\Delta$ 是由下述方式定义的普通范畴：
\begin{itemize}[leftmargin=2.2em,itemsep=0.1em]
\item 对象是二元组 $(X,E)$，其中 $X$ 是单纯集，$E\subset X_1$ 是 $X$ 中一族 $1$-单形，称为\textbf{带标记的 $1$-单形}集。
\item 从 $(X,E)$ 到 $(X',E')$ 的态射是 $\Set_\Delta$ 中的态射 $X\to X'$，使得 $E$ 被映入 $E'$。
\end{itemize}
我们称之为\textbf{标记单纯集}的范畴。
\end{definition}

\begin{construction}\label{L8-A.2}
有函子
\[
(-)^\flat\colon \Set_\Delta\to \Set^+_\Delta ,
\]
使得 $X^\flat$ 中的带标记 $1$-单形由 $X$ 中退化的 $1$-单形给出。

有函子
\[
(-)^\sharp\colon \Set_\Delta\to \Set^+_\Delta ,
\]
使得 $X$ 中每个 $1$-单形都是 $X^\sharp$ 中的带标记 $1$-单形。

有函子
\[
(-)^\natural\colon \mathrm{QCat}\to \Set^+_\Delta ,
\]
使得 $X^\natural$ 中的带标记 $1$-单形由 $X$ 中的同构给出。
\end{construction}

\begin{exercise}\label{L8-A.3}
设 $X$ 与 $Y$ 是 $\Set^+_\Delta$ 中的对象。证明函子
\[
\Hom_{\Set^+_\Delta}((-)^\flat\times Y,-)\colon \Set_\Delta^{\op}\to \Set
\]
可由 $\Set_\Delta$ 中的某个对象表示，我们把它记为 $\Fun^\flat(X,Y)$。

类似地，证明函子
\[
\Hom_{\Set^+_\Delta}((-)^\sharp\times Y,-)\colon \Set_\Delta^{\op}\to \Set
\]
可由 $\Set_\Delta$ 中的某个对象表示，我们把它记为 $\Fun^\sharp(X,Y)$。
\end{exercise}

\begin{exercise}\label{L8-A.4}
在 $\Set^+_\Delta$ 上构造典范模型结构，使得
\begin{itemize}[leftmargin=2.2em,itemsep=0.1em]
\item 任意态射都是上纤维化；
\item 双纤维性对象恰为 $X^\natural$，其中 $X\in\mathrm{QCat}$；
\item 双纤维性对象之间的弱等价由拟范畴之间的范畴等价给出；
\item 该模型结构与 $\Fun^\sharp(-,-)$ 给出的单纯充实相容。
\end{itemize}
\end{exercise}

\nitem\label{L8-A.5} \textbf{建议阅读}\quad HTT.3.1。
